% =====================================================================
\chapter{Probabilidad condicionada, independencia y teorema de Bayes}
\label{tema:4}
% =====================================================================

La probabilidad condicionada formaliza cómo cambia la incertidumbre al recibir información. De ella se deducen la regla de la cadena, el teorema de la probabilidad total y el teorema de Bayes, y con ella se define la independencia. El tema sigue las secciones §1.3 a §1.5 de Bertsekas y Tsitsiklis~\cite{bt} (clases~2 y~3 de MIT 6.041~\cite{mit6041f10}), el capítulo~2 de Blitzstein y Hwang~\cite{bh}, el capítulo~3 de Ross~\cite{ross} y el capítulo~17 de~\cite{mcs}, cuyo tratamiento del problema de Monty Hall y de las pruebas médicas (§17.1 y §17.4.2) se adapta aquí.

\section{Probabilidad condicionada}

\begin{definicion}[Probabilidad condicionada]
Sean $A$ y $B$ sucesos con $\PP(B)>0$. La \emph{probabilidad de $A$ condicionada a $B$} es
\[
\PP(A\mid B)=\frac{\PP(A\cap B)}{\PP(B)}.
\]
\end{definicion}

La idea es que, al saber que ha ocurrido $B$, los resultados fuera de $B$ quedan descartados y los de $B$ conservan sus probabilidades \emph{relativas}; dividir por $\PP(B)$ restablece la normalización. En el modelo de Laplace, $\PP(A\mid B)=\abs{A\cap B}/\abs{B}$: se cuenta dentro de $B$ como si $B$ fuera el nuevo espacio muestral.

\begin{proposicion}\label{prop:t4-esprob}
Si $\PP(B)>0$, la función $A\mapsto\PP(A\mid B)$ es una ley de probabilidad sobre $\Omega$. En consecuencia, cumple todas las propiedades de la \cref{prop:t3-props}; por ejemplo, $\PP(A^{c}\mid B)=1-\PP(A\mid B)$.
\end{proposicion}

\addtocounter{definicion}{1}%

\section{Regla de la cadena}

\begin{teorema}[Regla de la cadena o de la multiplicación]\label{teo:t4-cadena}
Si $A_1,\dots,A_n$ son sucesos con $\PP(A_1\cap\dots\cap A_{n-1})>0$, entonces
\[
\PP(A_1\cap A_2\cap\dots\cap A_n)=\PP(A_1)\,\PP(A_2\mid A_1)\,\PP(A_3\mid A_1\cap A_2)\cdots\PP(A_n\mid A_1\cap\dots\cap A_{n-1}).
\]
\end{teorema}

\addtocounter{definicion}{1}%

\section{Probabilidad total y teorema de Bayes}

\begin{definicion}
Una familia de sucesos $B_1,\dots,B_n$ es una \emph{partición} de $\Omega$ si son incompatibles dos a dos y $B_1\cup\dots\cup B_n=\Omega$.
\end{definicion}

\begin{teorema}[Probabilidad total]\label{teo:t4-total}
Si $B_1,\dots,B_n$ es una partición de $\Omega$ con $\PP(B_i)>0$ para todo $i$, entonces, para todo suceso $A$,
\[
\PP(A)=\sum_{i=1}^{n}\PP(A\mid B_i)\,\PP(B_i).
\]
\end{teorema}

\begin{teorema}[Bayes]\label{teo:t4-bayes}
En las condiciones del \cref{teo:t4-total}, si además $\PP(A)>0$,
\[
\PP(B_j\mid A)=\frac{\PP(A\mid B_j)\,\PP(B_j)}{\sum_{i=1}^{n}\PP(A\mid B_i)\,\PP(B_i)}.
\]
\end{teorema}

En la terminología bayesiana, $\PP(B_j)$ es la probabilidad \emph{a priori} de la hipótesis $B_j$, $\PP(A\mid B_j)$ es la \emph{verosimilitud} de la evidencia $A$ bajo esa hipótesis y $\PP(B_j\mid A)$ es la probabilidad \emph{a posteriori}. Con dos hipótesis $B$ y $B^{c}$, dividiendo las dos fórmulas de Bayes se obtiene la \emph{forma de razones} (\emph{odds}):
\[
\frac{\PP(B\mid A)}{\PP(B^{c}\mid A)}=\frac{\PP(B)}{\PP(B^{c})}\cdot\frac{\PP(A\mid B)}{\PP(A\mid B^{c})},
\]
«razón a posteriori = razón a priori $\times$ razón de verosimilitudes».

\addtocounter{definicion}{1}%

\addtocounter{definicion}{1}%

\section{Independencia}

\begin{definicion}[Independencia]
Dos sucesos $A$ y $B$ son \emph{independientes} si $\PP(A\cap B)=\PP(A)\PP(B)$. Si $\PP(B)>0$, esto equivale a $\PP(A\mid B)=\PP(A)$: saber que ha ocurrido $B$ no cambia la probabilidad de $A$. Los sucesos $A_1,\dots,A_n$ son \emph{independientes} (o mutuamente independientes) si para todo subconjunto de índices $S\subseteq[n]$ con $\abs{S}\ge2$,
\[
\PP\Bigl(\bigcap_{i\in S}A_i\Bigr)=\prod_{i\in S}\PP(A_i).
\]
Son \emph{independientes dos a dos} si la condición se cumple al menos para todos los $S$ con $\abs{S}=2$.
\end{definicion}

\begin{proposicion}\label{prop:t4-complementos}
Si $A$ y $B$ son independientes, también lo son $A$ y $B^{c}$, $A^{c}$ y $B$, y $A^{c}$ y $B^{c}$.
\end{proposicion}

\addtocounter{definicion}{1}%

\begin{observacion}[Independencia e incompatibilidad]
No deben confundirse. Si $A$ y $B$ son incompatibles con probabilidades positivas, \emph{no} son independientes: $\PP(A\cap B)=0<\PP(A)\PP(B)$; saber que ocurre $A$ garantiza que no ocurre $B$. La independencia es, además, una propiedad de la ley $\PP$, no de los sucesos como conjuntos: en la práctica suele ser una \emph{hipótesis} de modelización (lanzamientos sucesivos, ensayos repetidos) más que algo que se demuestre~\cite[§17.7.2]{mcs}.
\end{observacion}

Si $\PP(C)>0$, los sucesos $A$ y $B$ son \emph{condicionalmente independientes dado $C$} si $\PP(A\cap B\mid C)=\PP(A\mid C)\PP(B\mid C)$. Ni la independencia implica la independencia condicional ni al revés (\cref{eje:4-ajedrez}).

\addtocounter{definicion}{1}%

\addtocounter{definicion}{1}%

